% ============================================================
%   Appendix D  --  Transparent Use of AI
% ============================================================

\chapter{Transparent use of AI}\label{app:ai-usage}

\begin{figure}[H]
  \centering
  \includegraphics[width=0.6\textwidth]{imgs/ai-label_banner-assisted-by-ai.png}
  \caption{Disclosure banner indicating that parts of this project's development and
           documentation were produced with the assistance of AI tools, under continuous
           human supervision.}
  \label{fig:ai-banner}
\end{figure}

In line with the growing institutional expectation that academic work disclose the use
of generative AI tools, this appendix documents how, where, and why such tools were
employed during the development of \textit{CALA\_SafeRL}, the analysis of experimental
results, and the writing of this memoir.

\paragraph{Human oversight statement.}
It must be stated explicitly that \textbf{every output produced with the assistance of
an AI tool, whether source code, data analysis, or written text, was reviewed, tested,
and revised by the author before being incorporated into the project or this document}.
AI tools were used strictly in an assistive capacity, as accelerators for routine or
exploratory work, debugging aids, and proofreading support. At no point did an AI tool
act as an autonomous decision-maker with respect to the project's scientific
conclusions, the design of the safety shield, or the interpretation of results. Code
suggested by an assistant was run and inspected by the author; quantitative claims were
cross-checked against the underlying SQLite metric databases before being reported. This
principle of continuous human oversight is reiterated at the end of this appendix.

\section{Claude Code as a programming assistant}\label{app:ai-claude}

\textit{Claude Code}~\cite{anthropic2025claudecode} is an agentic command-line coding
assistant developed by Anthropic, used as the primary programming assistant during
implementation and experimentation. It operated directly inside the repository, with
access to the source tree, the test suite, and the per-run metrics database.

\paragraph{Debugging the wrapper chain.}
The layered \texttt{gymnasium} wrapper chain connecting the PPO agent to the CARLA
simulator (Chapter~\ref{ch:design}) is a frequent source of subtle bugs: the
agent interacts with \texttt{CarlaRewardShaper}, which wraps the active safety shield,
which in turn wraps \texttt{CarlaEnv}. An error in how one layer modifies the action,
observation, or reward can surface as a confusing symptom several layers away, for
example a mismatch between the policy's log-probabilities and the action actually
executed in the simulator (\texttt{info["executed\_action"]}).

The typical workflow was to describe a symptom in plain language, ask Claude Code to
trace the relevant code path across the wrapper boundaries, and then verify the proposed
explanation manually against the source code and the relevant unit test in
\texttt{tests/} before applying and re-testing a fix. This pattern, describe, trace,
verify, fix, was repeated for most of the non-trivial bugs in the shield credit-assignment
logic and in the continuous action-space handling.

\paragraph{Continuous action space and curriculum learning.}
Debugging the saturation of \texttt{log\_std} (the effective policy standard deviation
becoming pinned at its upper bound during early training) required correlating
per-update telemetry (\texttt{log\_std\_steering\_raw}, \texttt{log\_std\_throttle\_raw},
\texttt{log\_std\_saturated\_fraction}) with the entropy-coefficient schedule and the
\texttt{tanh}-squashed parameterisation in \texttt{ActorCritic.log\_std()}. Claude Code
was used to write throwaway scripts that queried \texttt{metrics.sqlite} and plotted
these quantities across runs, which sped up forming and discarding hypotheses. The
resulting decision, to replace the hard clamp with the smooth squashed parameterisation
described in Chapter~\ref{ch:design}, was the author's; the assistant's role was
limited to producing the exploratory tooling.

A similar use applied to the multi-phase curriculum in \texttt{train\_curriculum.py},
where each phase warm-starts from the previous phase's checkpoint. Claude Code helped
compare checkpoint metadata (\texttt{obs\_normalizer}, \texttt{ret\_rms},
\texttt{returns\_acc}) across phase boundaries, ruling out checkpoint-compatibility
issues and focusing the investigation on the hyperparameter changes themselves.

\paragraph{Large-scale analysis of PPO training metrics.}
Each training run produces per-episode and per-update SQLite databases
(Chapter~\ref{ch:design}), and many model variants were trained across shield
configurations, curriculum settings, and reward-shaping parameters. Claude Code drafted
SQL queries and analysis scripts to aggregate crash rate, off-road rate, mean reward,
and shielded fraction across runs, producing the comparison tables used in
Chapter~\ref{ch:experimentation}. As before, every aggregate figure reported in this
memoir was independently spot-checked by the author against the raw databases.

\section{Ollama with Gemma as a local, privacy-preserving assistant}\label{app:ai-ollama}

A second class of tasks benefited from running entirely on local hardware, using
\textit{Ollama} to host a \textit{Gemma}~\cite{team2024gemma} model.

\paragraph{Motivation.}
Two reasons motivated a local deployment. First, the per-run \texttt{metrics.sqlite}
databases contain project-specific experimental data, and keeping their analysis on the
local machine avoids any ambiguity about data retention by a third party. Second, the
\texttt{dataVisualizer.py} dashboard (Chapter~\ref{ch:design}) provides live
analysis while training is in progress, and a local model avoids the latency and
connectivity requirements of a remote API during long unattended runs.

\paragraph{Local analysis of training metrics.}
The dashboard optionally invokes the local Gemma model to narrate the Plotly charts
rendered from the live database, for example summarising whether a run's crash-rate
curve has plateaued or shows the rollback behaviour implemented by
\texttt{CurriculumManager}. The prompts include only numerical summaries already
computed locally (rolling crash/off-road rates, mean rewards, curriculum stage
transitions), not raw trajectory data. These summaries were treated as a first pass to
flag runs for closer inspection, never as a substitute for the quantitative analysis in
Chapter~\ref{ch:experimentation}.

\paragraph{Collision-type categorisation.}
A second, ongoing use of the local Gemma deployment is the categorisation of collision
events for the planned LLM-based reward-exposure work
(\texttt{tex/08\_llm\_reward\_exposure.tex}, in preparation). This part of the
project is still work in progress: the goal is to classify recorded collision episodes
into semantically meaningful categories, for example distinguishing a collision with an
oncoming vehicle during a forced lane change from one with a stationary obstacle near the
road edge, in order to inform the preference-based reward signal described in that
chapter. Running this locally allows the pipeline to be iterated on quickly without
per-call costs and keeps the per-collision context on the local machine. The categories
proposed by the model are a starting point for the preference-annotation pipeline and
remain subject to manual spot-checking, a validation step that is part of the work still
to be completed.

\section{Gemini as a writing and proofreading assistant}\label{app:ai-gemini}

\textit{Gemini}~\cite{geminiteam2025gemini}, developed by Google, was used during the
writing of this memoir as a structural and proofreading assistant, confined to the
presentation layer of the document.

Draft sections, written first by the author based directly on the implementation and on
the results in Chapter~\ref{ch:experimentation}, were passed to Gemini with requests
such as checking clarity and terminology consistency, or suggesting a clearer paragraph
structure without changing the technical content. Suggestions concerned sentence-level
clarity, consistent use of terms such as ``safety shield'', the ordering of paragraphs,
and minor grammatical corrections.

Every suggestion was reviewed by the author, who verified that no rewording altered the
technical meaning of a statement, introduced unsupported claims, or removed hedging
language around still-open questions, such as the work-in-progress status of the
collision-categorisation pipeline described in Section~\ref{app:ai-ollama}. Gemini
therefore improved readability and consistency, but did not influence the technical
content, argumentation, or conclusions of this memoir.

\section{Summary of human oversight}\label{app:ai-summary}

Three AI tools were used during this project, each in a distinct and bounded role:
Claude Code~\cite{anthropic2025claudecode} as a programming and data-analysis assistant
on the source code and metric databases; a locally-hosted Gemma
model~\cite{team2024gemma}, served through Ollama, for live metric narration and, in
ongoing work, collision categorisation; and Gemini~\cite{geminiteam2025gemini} as a
proofreading and structural-editing assistant for this memoir. In every case the role
was assistive, accelerating routine or exploratory tasks without making design
decisions, determining experimental conclusions, or producing final text or code without
human review. \textbf{All code, all quantitative results, and all written content in
this project and in this memoir were ultimately validated, tested, and approved by the
author}, who takes full responsibility for their correctness and content.
